\section{Gestión de la base de datos}
\label{sec:gestion_base_datos}

Los registros generados durante los ensayos de impulsos atmosféricos de alta tensión requieren un sistema de almacenamiento fiable y trazable. Para ello, el módulo de persistencia denominado \textbf{FileManager}, maneja de forma centralizada la persistencia de datos en disco, protegiendo la integridad de las formas de onda originales y los resultados analíticos calculados, evitando pérdidas de información ante fallos, a la vez que optimiza el uso de espacio en disco y reduce los tiempos de lectura/escritura.

\subsection{Flujo de almacenamiento}
\label{subsec:flujo_almacenamiento}

El flujo de guardado inicia cuando el operador acciona el comando en la interfaz gráfica, entonces, el Presentador ejecuta una secuencia de validación:
\begin{enumerate}
    \item \textbf{Verificación de sesión activa}: Comprueba la existencia de un proyecto nuevo o cargado desde memoria.
    \item \textbf{Validación sintáctica de campos}: Certifica que los campos obligatorios (datos del ensayo y condiciones ambientales) hayan sido cargados, cumpliendo con los filtros sintácticos aplicados mediante expresiones regulares.
    \item \textbf{Integridad del buffer}: Evalúa que el buffer temporal de adquisición contenga ondas digitalizadas pendientes de almacenamiento.
\end{enumerate}

Si alguna de las comprobaciones falla, el flujo se interrumpe y se emite una notificación de error en la interfaz de gráfica. Por el contrario, si se superan los filtros de validación, el Presentador clasifica el impulso (ensayo o referencia), actualiza el contador de onda correspondiente, y genera una marca temporal unívoca que asegura la trazabilidad del impulso.

Posteriormente, el módulo de almacenamiento, graba la información en disco, implementando mecanismos de tolerancia a fallos (como pérdidas de alimentación eléctrica o excepciones en tiempo de ejecución) para mitigar riesgos de corrupción de datos:

\begin{enumerate}
    \item \textbf{Apertura segura}: El archivo HDF5 se manipula en \emph{modo adición} usando un manejador de contexto, que asegura la
    escritura segura frente a fallos y la correcta liberación del descriptor de archivo del sistema operativo ante cierres inesperados, evitando la corrupción de la base de datos.
    \item \textbf{Organización jerárquica}: Se crea un subgrupo individual para cada registro (por ej. \texttt{/Onda\_01} o \texttt{/Referencia\_01}). Las condiciones ambientales, la fecha y los parámetros calculados, se inyectan como atributos de este grupo específico, garantizando que cada onda digitalizada esté inseparablemente asociada a sus parámetros de ensayo.
    \item \textbf{Resguardo de seguridad redundante}: En paralelo a la base de datos, crea copias de seguridad independientes de cada onda, en formato de texto plano con valores separados por comas (\textbf{CSV}), que contienen el registro bruto escalado verticalmente a su magnitud real. Lo que brinda un sistema de respaldo redundante de los datos originales, desacoplado del formato principal.
\end{enumerate}

\subsubsection{Exportación de resultados}
\label{subsec:estructura_base_datos}

Para exportar los resultados a hojas de cálculo que faciliten su manipulación en programas comerciales para la elaboración de informes, el módulo permite la extracción del historial completo de los impulsos aplicados al elemento bajo ensayo, hacia hojas de cálculo utilizando los formatos de extensión \textbf{CSV} o \textbf{XLSX}.

Aunque previamente convierte la escala física de los resultados: las bases de tiempo se convierten de segundos (\unit{\second}) a microsegundos (\unit{\micro\second}) y las tensiones de volts (\unit{\volt}) a kilovolts (\unit{\kilo\volt}).


\subsection{Sistema de ficheros en disco}
\label{subsec:estructura_ficheros_disco}

Una vez que el operador haya seleccionado un directorio raíz para el proyecto, la creación de carpetas se automatiza aplicando el árbol de directorios de la \autoref{fig:estructura_carpetas}, asignando espacios físicos en disco para la información.

\begin{figure}[H]
    \centering
    \begin{forest}
        for tree={
            folder,
            grow'=0,
            font=\sffamily\small,
            inner sep=3pt,
            s sep=3pt,        % Separación vertical ajustada
            draw=black!20,    % Color del borde del nodo (opcional)
            fill=black!5,     % Color de fondo del nodo (opcional)
            rounded corners   % Bordes redondeados (opcional)
        }
        [\textbf{Directorio raíz (Carpeta item-cliente)}
            [\textbf{01 Respaldo}: \emph{Carpeta de respaldo en texto plano (\textbf{CSV})}
            [{\textbf{\texttt{Onda\_nn\_AAAAMMDD\_HHMMSS\_V.csv}}: \emph{Archivo respaldo de onda de tensión}}]
            [{\textbf{\texttt{Onda\_nn\_AAAAMMDD\_HHMMSS\_A.csv}}: \emph{Archivo respaldo de onda de corriente}}]
            ]
            [\textbf{02 Analisis de datos}: \emph{Carpeta de datos procesados}
            [{\textbf{\texttt{item-cliente.h5}}: \emph{Base de datos en formato HDF5}}]
            ]
            [\textbf{03 Resultados}: \emph{Carpeta de resultados exportados (hoja de cálculo e imágenes)}
            [{\textbf{\texttt{item-cliente - Resultados.xlsx}}: \emph{Archivo de resultados del ensayo}}]
            [{\textbf{\texttt{item-cliente - Resultados.cvs}}: \emph{Archivo de resultados del ensayo}}]
            [{\textbf{\texttt{ondas.png}}: \emph{Imagen exportada manualmente desde el visualizador de ondas}}]
            ]
        ]
    \end{forest}
    \caption{Esctructura de directorio de archivos.}
    \label{fig:estructura_carpetas}
\end{figure}
Donde:
\begin{itemize}
    \item \texttt{nn}: Indica el número de onda con rango 00-99.
    \item \texttt{AAAAMMDD}: Indica la fecha en formato Año-Mes-Día.
    \item \texttt{HHMMSS}: Indica el horario en formato Hora-Minutos-Segundos.
\end{itemize}

\subsection{Estructura de la base de datos}
\label{subsec:estructura_base_datos}

La información se almacena de manera estructurada y comprimida, en único archivo con formato \textbf{HDF5}, que actúa como una base datos independiente para cada elemento bajo ensayo.

Este formato organiza internamente el archivo como un sistema de ficheros, distribuyendo los datos en grupos que actúan como carpetas, aplicando la estructura de la \autoref{fig:estructura_base_datos}.

\begin{figure}[H]
    \centering
    \begin{forest}
        for tree={
            folder,
            grow'=0,
            font=\sffamily\small,
            inner sep=3pt,
            s sep=3pt,        % Separación vertical ajustada
            draw=black!20,    % Color del borde del nodo (opcional)
            fill=black!5,     % Color de fondo del nodo (opcional)
            rounded corners   % Bordes redondeados (opcional)
        }
        [\textbf{Grupo raíz (Archivo .h5)}
            [\textbf{Atributos globales}: \emph{Datos comunes a todos los registros}
            [{\textbf{\texttt{Item}}: \emph{Identificador del objeto a ensayar}}]
            [{\textbf{\texttt{Cliente}}: \emph{Nombre de la entidad solicitante}}]
            [{\textbf{\texttt{Atenuaciones}}: \emph{Factores de atenuación de CH1 y CH2}}]
            ]
            [{\textbf{Grupos de ondas} \emph{Impulso capturado} (\texttt{Onda\_n} / \texttt{Referencia\_n})}
            [\textbf{Atributos particulares}: \emph{Metadatos propios de cada impulso}
                [{\textbf{\texttt{Fecha}}: \emph{Marca de tiempo de adquisición}}]
                [{\textbf{\texttt{Temperatura}}: \emph{Parámetro ambiental}}]
                [{\textbf{\texttt{Humedad}}: \emph{Parámetro ambiental}}]
                [{\textbf{\texttt{Presión}}: \emph{Parámetro ambiental}}]
                [{\textbf{\texttt{Polaridad}}: \emph{Sentido eléctrico del impulso} (+/-)}]
                [{\textbf{\texttt{Valor\_pico}}: \emph{Tensión de cresta}}]
                [{\textbf{\texttt{T1}}: \emph{Tiempo de frente}}]
                [{\textbf{\texttt{T2}}: \emph{Tiempo de cola}}]
                [{\textbf{\texttt{OS}}: \emph{Sobreimpulso porcentual}}]
            ]
            [{\textbf{Grupo de Tiempos} (\texttt{/Time}) [\unit{\second}]}
                [{\textbf{\texttt{raw}}: \emph{Vector bruto a apartir de} $\Delta t$}]
                [{\textbf{\texttt{alineado}}: \emph{Vector alineado con el frente del impulso de referencia}}]
            ]
            [{\textbf{Grupo de CH1} (\texttt{/CH1\_Voltage}) [\unit{\volt}]}
                [{\textbf{\texttt{raw}}: \emph{Vector de tensión bruta, magnitud real}}]
                [{\textbf{\texttt{ensayo}}: \emph{Vector de tensión de ensayo IEC 60060-1}}]
                [{\textbf{\texttt{normalizado}}: \emph{Vector normalizado} ($V_{\text{pico}} = \num{1.0}$)}]
            ]
            [{\textbf{Grupo de CH2} (\texttt{/CH2\_Current}) [\unit{\ampere}]}
                [{\textbf{\texttt{raw}}: \emph{Vector de corriente bruta, magnitud real}}]
                [{\textbf{\texttt{ensayo}}: \emph{Copia del bruto (simetría de procesamiento con CH1)}}]
                [{\textbf{\texttt{normalizado}}: \emph{Vector normalizado} ($I_{\text{pico}} = \num{1.0}$)}]
            ]
            ]
        ]
        \end{forest}
    \caption{Estructura jerárquica de la base de datos.}
    \label{fig:estructura_base_datos}
\end{figure}